\section{The renormalization group transformation}\label{sec:rg-step}

This part carries out one step of Ba\l aban's renormalization group for pure Yang--Mills theory in four dimensions in the small-field region \cite{BalYM4I,BalYM4II}; the three-dimensional version is \cite{BalUV3}. Here we derive an exact representation of one step as a Gaussian integral around the background field, and we prove the linear estimates on which the analysis of that integral rests:
\begin{itemize}
\item the coercivity and decay of the fluctuation covariance;
\item the analyticity of the coarse effective action and the decay of its derivatives;
\item the structure of the background-dependent normalization.
\end{itemize}
The polymer expansion of the interaction, its renormalization and the iteration are the subject of Parts~VIII--X.

\paragraph{What one step does.} The step integrates out the fluctuation of $V$ around the base configuration $W=U'^{(k)}$, the $k$-step average of the background field $U'$ of $V'$ (Section~\ref{ss:31.3}). After Ba\l aban's gauge fixing relative to $W$ and the elimination of the constraint on the middle bond $b_0(c)$ of each coarse bond $c$, it is an integral over fields $Z$ on the set $\mathfrak F$ of the other bonds (Section~\ref{sec:fluctuations}). For a density of the form \eqref{eq:rho-form} the integrand contains the factor $e^{-g_k^{-2}F(Z)}$, with $F(Z)=S(U')+\frac12\Delta_k(Z,Z)+\mathcal V(Z)+G_3(Z)$, where $\mathcal V$ and $G_3$ are of third order (Lemma~\ref{lem:step-quadratic}). Section~\ref{sec:coarse-action} studies the coarse action that enters $\Delta_k$, and Section~\ref{sec:covariance} proves that $\Delta_k$ is coercive and that $C_k=\Delta_k^{-1}$ decays (Corollary~\ref{cor:covariance}). For $k\ge1$, Theorem~\ref{thm:step} then writes the step as an explicit prefactor times a Gaussian integral with covariance $C_k$, and Theorem~\ref{thm:K} collects the results of this part. The \hyperref[intro:VII]{introduction to Part~VII} lists the stages of the step with their results.

\paragraph{Conventions for Part VII.} We keep the conventions of Parts~IV--VI. In particular $d=4$ and $L\ge3$ is odd (Section~\ref{sec:setup}). In addition $G$ is semisimple. Semisimplicity is not used in this part; it is needed for the renormalization in the next one. We work in the \emph{units of step $k$}: the unit lattice is $\T=\T^{(k)}$, the fine lattice is $\Tf$ with $\eta=L^{-k}$, and the next coarse lattice is $\T'=\T^{(k+1)}$, of spacing $L$. Its blocks are $B'(y)=B_{k+1}(y)$, with $L^d$ points of $\T$, and the one-step averaging is $\mathfrak a'=\mathfrak a_{k+1}$. The torus $\T$ has side $Ln$, where $n$ is the side of $\T'$ in its own units. We assume $n\ge n_{\delta=1}$, with $n_\delta$ from Proposition~\ref{prop:sharp-curvature}, and $n\ge80M'_L$, with $M'_L$ from Theorem~\ref{thm:step-coercive}. Fields on $\T$ carry the unweighted norm $\|Z\|^2=\sum_b|Z(b)|^2$.

\subsection{The model and the transformation}\label{ss:31.1}

Let $\epsilon=L^{-N}$ be the finest lattice spacing and $g$ the bare coupling. Here $N$ is the total number of scales (Section~\ref{subsec:tf-scales}); it is not the integer $N=L^k=\eta^{-1}$ of Section~\ref{sec:setup}, which appears in Section~\ref{ss:31.2}. The initial density on the $\epsilon$-lattice is $\rho_0(U)=\exp(-g^{-2}S_{\rm lat}(U)-E_0)$, where $S_{\rm lat}(U)=\sum_p(\mathsf N-\operatorname{Re}\operatorname{tr}U(\partial p))$. In the units of step $k$ the $\epsilon$-lattice has spacing $\eta=L^{-k}$, and $S_{\rm lat}=S$, the action of Section~\ref{sec:action}: in four dimensions the Wilson action carries no power of the lattice spacing. So $\rho_0=\exp(-g^{-2}S-E_0)$ in the units of every step, and a change of units alone does not change the coupling.

The averaging $\mathfrak a'$ uses group averages and logarithms, so it is defined only on part of the configuration space. Let $\operatorname{dom}\mathfrak a'$ be the set of real configurations $V$ on $\T$ for which Definition~\ref{def:one-step} defines $\mathfrak a'(V)$: for every positively oriented bond $c$ of $\T'$ and every $x\in B'(c_-)$, the averaged contour variables in $\Lambda_{c,x}(V)$ are defined, i.e.\ their families lie in $\cD_{r_M}$ (Definition~\ref{def:avg-contour}), and $|\Lambda_{c,x}(V)-1|\le c_G$. For real $V$ these averaged contour variables are real (Lemma~\ref{lem:group-average}(b)), so the logarithms $\ell(\Lambda_{c,x}(V))$ lie in $\mathfrak g$, and $\mathfrak a'(V)$ is defined and real, also on the negatively oriented bonds (Lemma~\ref{lem:reversal}). Outside this set an average or a logarithm may be undefined, or a logarithm may leave $\mathfrak g$. The set $\operatorname{dom}\mathfrak a'$ is a relatively closed subset of the open set on which the averaged contour variables are defined, and $\mathfrak a'$ is continuous on it (Lemma~\ref{lem:group-average}(a)). It is invariant under real gauge transformations and under the lattice symmetries $\mathfrak S_{k+1}$ (Lemma~\ref{lem:one-step-basic}(b),(d)). It contains $\mathfrak U_k(b)$ in its interior whenever $8bL^2\le r_M$: then, by Lemma~\ref{lem:stokes-avg}(i),(iii) and the proof of Lemma~\ref{lem:one-step-basic}(a), with $\delta_{k+1}=L$, the families lie in $\cD_{r_M/2}$ and $|\Lambda_{c,x}(V)-1|\le2r_M\le\frac12c_G$.

For an integrable density $\rho$ on configurations $V$ on $\T$ that vanishes outside $\operatorname{dom}\mathfrak a'$, let $T\rho$ be the image of the measure $\rho(V)\prod_bdV(b)$ under $\mathfrak a'$:
\[
\int\varphi\,d(T\rho)=\int\prod_bdV(b)\ \varphi\big(\mathfrak a'(V)\big)\,\rho(V)
\]
for bounded Borel functions $\varphi$ of configurations on $\T'$. Here $b$ runs over the positively oriented bonds of $\T$, and $dV(b)$ is the normalized Haar measure. With $\varphi=1$, the total mass of $T\rho$ is $\int\prod_bdV(b)\,\rho(V)$. We write
\begin{equation}\label{eq:rg-map}
(T\rho)(V')=\int\prod_{b}dV(b)\ \prod_{c}\delta\big(\mathfrak a'(V)(c)\,V'(c)^{-1}\big)\ \rho(V),
\end{equation}
where $c$ runs over the positively oriented bonds of $\T'$ and $\delta$ is the $\delta$-function of $G$ at $1$ with respect to the Haar measure. This formula is a notation for the image measure. The left side denotes a density of $T\rho$ with respect to $\prod_cdV'(c)$, and it is used only where such a density exists; the remark at the end of this subsection says in what sense the formula holds, and when the density exists. Passing to the units of step $k+1$ relabels the configurations, without any Jacobian since they are group valued, and leaves $S$ unchanged, since $L^{4-d}=1$.

\paragraph{The couplings.} Unlike in three dimensions, the coupling of step $k+1$ is therefore not produced by the change of units. It is produced only by the integration over the fluctuation field. The effective action of step $k+1$ contains a term proportional to $S$, with a coefficient $\beta$ of order one for fixed $L$, which Ba\l aban reads off from the vacuum polarization \cite[(1.20)--(1.22)]{BalYM4I}, and the inverse square of the coupling of step $k+1$ is that of step $k$ corrected by $\beta$. The sequence of couplings is defined together with the generations of the flow, in Part~X: $\beta_{k+1}$ is read off from the period density of the new generation, and $g_{k+1}^{-2}=g_k^{-2}-\beta_{k+1}$ (Proposition~\ref{prop:period-iso}, Theorems~\ref{thm:gen-step} and~\ref{thm:sf-flow}). Part~XI computes the one-loop value of $\beta$ and compares it with the one-loop parts of the coefficients of the flow (Section~\ref{sec:ol-uv}). In this part $g_k$ is a parameter: the coupling of the density at step $k$, with $g_0=g$. The statements about couplings and thresholds below hold for every such parameter that satisfies their hypotheses; along the flow these hypotheses hold at every step (Theorem~\ref{thm:sf-flow}(e)).

If $\rho$ vanishes outside $\operatorname{dom}\mathfrak a'$ and $\rho(V^u)=\rho(V)$ for all real gauge transformations $u$ on $\T$, then $T\rho$ is gauge invariant on $\T'$: it is invariant under $V'\mapsto V'^{u'}$ for every real gauge transformation $u'$ on $\T'$. This follows from $\mathfrak a'(V^u)=\mathfrak a'(V)^{u|_{\T'}}$ for real $u$ of any size (Lemma~\ref{lem:one-step-basic}(b)), applied to an extension $u$ of $u'$, from the invariance of $\operatorname{dom}\mathfrak a'$, and from the invariance of the Haar measure under $V\mapsto V^u$. All densities below are gauge invariant. The map $T$ also commutes with the lattice symmetries of $\T'$ (Proposition~\ref{prop:symmetry}(a)).

\paragraph{The small-field step.} In the small-field region the density at step $k$ has, by assumption, the form
\begin{equation}\label{eq:rho-form}
\rho_k(V)=\exp\Big(-\frac1{g_k^2}S\big(U_k(V)\big)+\mathfrak E_k(V)\Big)\qquad\text{for }V\in\mathfrak U_k(\tilde h),
\end{equation}
with $U_k(V)$ the background field of Theorem~\ref{thm:balaban}, a gauge-invariant function $\mathfrak E_k$, measurable and bounded on $\mathfrak U_k(\tilde h)$, and a threshold $\tilde h\le a_1$. The renormalization group step studied here is the \emph{small-field step}
\[
T_{\rm sf}(\rho):=T^{\rm gf}\big(\chi^{\rm w}\chi_\xi\,\rho\big)
\]
(Definitions~\ref{def:gf-step}, \ref{def:chi-w} and~\ref{def:small-field-step}). Here $T^{\rm gf}$ is the map \eqref{eq:rg-map} with Ba\l aban's gauge-fixing function in the integrand, $\chi^{\rm w}$ restricts the fluctuation field around the background of the step to $|\cdot|\le w$, and $\chi_\xi$ imposes the small-field condition on the white noise, as in \cite[(206)]{DimockI}. The gauge-fixing function and the cutoffs depend on $V'$ through the base configuration and the reference of the step. In $T^{\rm gf}$ they are taken at $V'=\mathfrak a'(V)$. The gauge-fixing function is set to $0$ unless $V\in\operatorname{dom}\mathfrak a'$ and $\mathfrak a'(V)$ is as in Lemma~\ref{lem:composition}, so that the base is defined, and the cutoffs are set to $0$ unless $\mathfrak a'(V)$ satisfies the hypotheses of Proposition~\ref{prop:exact-step}. That proposition shows that, for such $V'$, $\chi^{\rm w}$ vanishes outside Ba\l aban's set $\mathfrak M$ of regular configurations (Section~\ref{sec:fluctuations}), which lies in the interior of $\operatorname{dom}\mathfrak a'$, so that the restriction to $\operatorname{dom}\mathfrak a'$ removes nothing, and that $\mathfrak a'$ is a submersion wherever $\chi^{\rm w}=1$. Regular configurations in this sense are not the $(M,\eps)$-regular backgrounds of Definition~\ref{def:regular}. These functions also depend on the choice of $U'$ in the minimal orbit of $V'$. We fix a Borel choice; it exists by the remark at the end of this subsection. The density at each $V'$ does not depend on the choice, by Lemma~\ref{lem:gf-step-invariance}(i). The cutoffs act bondwise on the gauge-fixed variables, as Ba\l aban's cutoff does \cite[(2.9), (2.10)]{BalYM4I}. In particular $\chi^{\rm w}$ is not invariant under the gauge transformations of $\T$ that are trivial on $\T'$, so $T_{\rm sf}$ is a modified model: it is not of the form $T(\chi\rho)$ with an invariant cutoff $\chi$. On integrands that are invariant under these gauge transformations and supported on regular configurations (Section~\ref{sec:averaging}), $T^{\rm gf}=T$ (Lemma~\ref{lem:exp-gauge}(a)). The step is the analogue of \cite[(172), (206)]{DimockI}; Ba\l aban's version is \cite[(7)--(13)]{BalUV3} in three dimensions and \cite[(2.9), (2.10)]{BalYM4I} in four. The large-field complement is not treated here.

\begin{remark*}[The image measure and the Borel choice]
\emph{The formula \eqref{eq:rg-map}.} If $\delta$ is replaced by a continuous function $\delta_n\ge0$ on $G$ with $\int\delta_n=1$, supported in a neighbourhood of $1$ that shrinks as $n\to\infty$, the right side is a continuous density, and the measures with these densities converge weakly to $T\rho$: for continuous $\varphi$ this follows from Fubini's theorem, the invariance of the Haar measure and the uniform continuity of $\varphi$. A sufficient condition for the density $(T\rho)(V')$ to exist is that $\rho$ vanishes outside an open subset of $\operatorname{dom}\mathfrak a'$ on which $\mathfrak a'$ is a submersion, i.e.\ its derivative has rank $\dim\mathfrak g$ times the number of positively oriented bonds of $\T'$. Then, near each point of this set, $\mathfrak a'$ is a coordinate projection in suitable analytic coordinates, and for almost every $V'$ the density at $V'$ is the integral of $\rho$ over the fibre $\{\mathfrak a'(V)=V'\}$ with respect to the measure that the two Haar measures induce on it. Integrability of $\rho$ alone does not make this fibre integral finite at every $V'$. If $\rho$ is a bounded Borel function and vanishes outside a compact subset of the open set, the fibre integral converges absolutely at every $V'$; it is then taken as the density at every $V'$, and the usual rules for $\delta$-functions compute it. In that case $\int\prod_cdV'(c)\,(T\rho)(V')=\int\prod_bdV(b)\,\rho(V)$. The density is used for the small-field step. There the integrand is bounded, and the cutoff confines it to a compact subset of the region of Lemma~\ref{lem:eliminate}, which gives the submersion, and Proposition~\ref{prop:exact-step} computes the density. Part~XI uses it also for integrands supported in the gauge orbits of that region (proof of Lemma~\ref{lem:ol-one-step}), where the same argument applies.

\emph{The Borel choice of $U'$.} A Borel choice exists: near each $V'$, Corollary~\ref{cor:analytic} gives a choice that depends continuously on $V'$, and finitely many such neighbourhoods cover the compact set of the $V'$ as in Lemma~\ref{lem:composition}. For such a choice the gauge-fixing function and the cutoffs are built from continuous maps ($\mathfrak a'$, the base, the averaged contour variables, $\ell$ and $C_k^{-1/2}$, which depends continuously on the base by Corollary~\ref{cor:covariance}), continuous functions and indicator functions of closed sets. So they are Borel functions of $V$.
\end{remark*}

\subsection{Curvature of background fields}\label{ss:31.2}

Theorem~\ref{thm:balaban} bounds the curvature of $U_k(V)$ by $B_3\eps_1$, where $B_3$ depends on $L$. The small-field thresholds of consecutive steps are compatible only if the curvature of the background is bounded with a constant that does \emph{not} grow with $L$. That is proved here.

The constants of the pointwise bounds of Parts~II and~III must therefore be independent of $L$ and of $M$; note that $M\ge5L$ in Part~VI. We make this explicit. Ba\l aban's averaging pair enters only through one parameter.

\begin{definition}[Parameter of an averaging pair]\label{def:avg-parameter}
Let $U$ be a real $(1,\eps)$-regular background and $(Q',\bQ)$ an averaging pair with properties (i), (iii) and (iv) of Remark~\ref{rem:general-averaging}. Let $\eps_Q\ge0$. The pair has \emph{parameter $\eps_Q$ at $U$} if the following holds for every $y\in\T$ and every real gauge transformation $u$ in which (B) holds with parameter $\eps$ on $\Lam_4(c_y)$. Here (B) on $\Lam_4(c_y)$ refers to the bonds and differences that involve only points of $\Lam_4(c_y)$, and the H\"older quotient on a cube $\Lam_3(z'')\subset\Lam_4(c_y)$ is taken over the pairs of differences that involve only points of $\Lam_3(z'')$. If $U$ is $(M,\eps)$-regular, every gauge of Definition~\ref{def:regular} on a cube containing $\Lam_4(c_y)$ satisfies (B) with parameter $\eps$ on $\Lam_4(c_y)$ in this sense. Computed in the gauge $u$, on the blocks and tubes contained in $\Lam_4(c_y)$:
\begin{enumerate}[label=(\alph*)]
\item $Q'(U^u)-Q_0$ and $\bQ(U^u)-\bQ$ satisfy the bounds of Lemma~\ref{lem:cov-avg}(b) with $C\eps$ replaced by $\eps_Q$: block norms, the bounds from sup norms to values, and the same for the transposes;
\item the defect $E$ satisfies the bounds of Lemma~\ref{lem:cov-avg}(c) with $C\eps$ replaced by $\eps_Q$, and its transpose satisfies $|E^T\omega|_\infty\le\eps_Q|\omega|_\infty$.
\end{enumerate}
\end{definition}

By Lemma~\ref{lem:cov-avg}, the model pair $(Q'(U),\bQ(U))$ of Section~\ref{sec:background} has parameter $C_{\rm mod}\eps$, with an absolute constant $C_{\rm mod}$.

\begin{lemma}[Monotonicity of regularity]\label{lem:regular-monotone}
If $U$ is $(M,\eps)$-regular and $1\le M'\le M$, then $U$ is $(M',\eps)$-regular.
\end{lemma}

\begin{proof}\hasdetail{lem:regular-monotone}
The cube $\{x:|x-z|_\infty\le4M'+2\}$ lies in $\{x:|x-z|_\infty\le4M+2\}$. Restrict the gauge transformation of Definition~\ref{def:regular} to it. The bonds with both endpoints in the smaller cube have both endpoints in the larger one, and every cube $\Lam_3(z')$ in the smaller cube lies in the larger one. So the three bounds of Definition~\ref{def:regular} persist.
\end{proof}

\begin{lemma}[Uniformity of the constants of Parts~II--III]\label{lem:uniformity}
Fix $a>0$ and $\alpha\in(0,1)$. Let $M_0$ be the constant of Part~III (Theorem~\ref{thm:reg-coercive}). There are $\eps_{\rm U}$, $\theta_Q$, $C_H$, $\gamma_H>0$, depending only on $a$, $\alpha$ and $\dim\mathfrak g$, with the following property. Let $k\ge1$ and let the side of $\T$ be at least $8M_0$. Let $M\ge M_0$, let $U_0$ be real and $(M,\eps)$-regular with $\eps\le\eps_{\rm U}$, and let $(Q',\bQ)$ be an averaging pair with properties (i)--(iv) of Remark~\ref{rem:general-averaging} and parameter $\eps_Q\le\theta_Q$ at $U_0$. Then:
\begin{enumerate}[label=(\alph*)]
\item Theorem~\ref{thm:reg-scalar} (with $X=\Tf$), Theorems~\ref{thm:reg-coercive} and~\ref{thm:reg-gauge} (with $X=\Tf$), and Theorem~\ref{thm:reg-min} hold at $U_0$ with this pair, in the form stated there: existence, covariance, identities, coercivity, decay and the pointwise bounds in local gauges, with no comparison with the operators at $U=1$. Their constants $c,C,\gamma,c_*$ depend only on $a$, $\alpha$, $\dim\mathfrak g$: neither on $M$, nor on $N=L^k$, $n$, $k$, nor on the pair beyond $\eps_Q\le\theta_Q$.
\item Let $\bH=\bH(U_0)$ be the minimizer map of Theorem~\ref{thm:reg-min} (with $W=0$) and $\nabla$ the covariant difference at $U_0$. For every unit-lattice field $B$,
\begin{equation}\label{eq:H-pointwise}
|\bH B|_{B(y)}+|\nabla\bH B|_{B(y)}\le C_H\sum_{y'}e^{-\gamma_H|y-y'|}|B|_{y'} .
\end{equation}
\item For Ba\l aban's pair $(Q'_k(U_0),\bQ_k(U_0))$ we have $E=0$, and there are $c_Q,C_Q>0$ such that the pair has parameter $\eps_Q\le C_Q\eps$ for $\eps\le c_Q$. Here $c_Q$, $C_Q$ may depend on $d$, $L$, $G$, $a$, $\alpha$ and $M$ (the conventions of Part~IV). Hence (a) and (b) hold for this pair whenever $\eps\le\eps_{\rm B}=\min(\eps_{\rm U},\theta_Q/C_Q,c_Q)$. In particular \eqref{eq:H-pointwise} holds with $C_H$, $\gamma_H$ independent of $L$ and $M$, for $\eps$ small depending on $L$ and $M$.
\end{enumerate}
If the side of $\T$ is less than $8M_0$, (a)--(c) hold as well, with the same constants.
\end{lemma}

\begin{proof}\hasdetail{lem:uniformity}
\emph{Step 0: reduction to $M=M_0$.} For $X=\Tf$ every $j$ lies in $J(\Tf)$, so $m_{\Tf}=0$. Hence the torus operators $\fA'_{\Tf}$, $G'$, $\Gamma'$, $P_\cV$, $R$, $S$, $\Pi$, $\fA_\pi=\fA^{\Tf}$, $\bG$, $\boldsymbol\Gamma_\pi$, $\bH$, $\bC$ are defined without reference to $M$. By Lemma~\ref{lem:regular-monotone}, $U_0$ is $(M_0,\eps)$-regular. Apply every result of Part~III with $M=M_0$; their hypotheses are ``$M\ge M_0$ and $(M,\eps)$-regular''. Their constants then depend on $M_0$, which depends only on $a,\alpha,\dim\mathfrak g$. Coercivity and decay are properties of the fixed torus operators, so proving them at scale $M_0$ proves them for $U_0$. This proves the independence of $M$. The scale $M_0$ also enters through the partition of unity (Lemma~\ref{lem:partition}), which needs the side of $\T$ to be at least $8M_0$. If the side is smaller, the partition at scale $M_0$ is the trivial one of Section~\ref{ss:17.2}. The cube of Definition~\ref{def:regular} at scale $M_0$ then covers $\Tf$, so one real gauge gives the bounds of that definition on all of $\Tf$. The proof of Theorem~\ref{thm:reg-coercive} treats this case, and the other results in (a) do not use the partition. In Steps~1--4 this gauge serves as the local gauge of every cube, and (a)--(c) follow with the same constants.

\emph{Step 1: Part~I inputs.} The model pair at $U=1$ is $(Q_0,\bQ)$ on centred blocks of side $1$ with $N^d$ fine points. The following constants are used, and none depends on the odd integer $N=L^k\ge3$ (Section~\ref{sec:setup}: $L$ enters only through $N$, and $k\ge1$). In particular the margins of the trial fields of Theorem~\ref{thm:reg-scalar}(c) and Lemma~\ref{lem:trial-gauge}, built from products of $\cos^2$, hold for every odd $N\ge3$ with the same constants.
\begin{itemize}
\item the block Poincar\'e constants (Lemma~\ref{lem:block-poincare}, Corollary~\ref{cor:poincare-all}), which are absolute;
\item the gauge-field coercivity constant $c_g$ (Theorem~\ref{thm:gauge-coercive}), depending on $a$, on the absolute constants $C_P$, $C_W$ (Lemma~\ref{lem:Q-divergence}) and $C_B$ (Lemma~\ref{lem:biharmonic}), and on the bound in Proposition~\ref{prop:R}(e);
\item the decay of $P_\cV(1)$ and $dP_\cV(1)d^*$ (Proposition~\ref{prop:R}), depending on $a$;
\item the Combes--Thomas constants (Lemmas~\ref{lem:CT-bounded}, \ref{lem:neumann}, \ref{lem:CT-form}), explicit in the coercivity and conjugation constants;
\item the interior regularity constant $C_{\rm reg}$ (Theorem~\ref{thm:local-regularity}), depending on $\alpha$, for all $N\ge1$.
\end{itemize}
All these are statements on the fine lattice with unit blocks. Distances, decay rates and the norms $|\cdot|_{B(y)}$ are measured in unit-lattice units. The supports of the averaging kernels, $B(y)$ and $S_c=B(c_-)\cup B(c_+)$, consist of a fixed number of unit blocks, whatever $L$ is.

\emph{Step 2: where the averaging pair enters.} We list every use of the pair in the proofs of the results in (a). In each, the model bound $C\eps$ is replaced by $\eps_Q$, and the conclusion holds once $\eps_Q\le\theta_Q$, with $\theta_Q$ a threshold of the same kind as $\eps_{\rm U}$.
\begin{enumerate}[label=(A\arabic*)]
\item Theorem~\ref{thm:reg-scalar}(a): in the local gauge, $|(Q'\lambda)(y)-(Q_0\tilde\lambda)(y)|\le\eps_Q\|\lambda\|_{B(y)}$, by Definition~\ref{def:avg-parameter}(a) and the locality (iv). The error term becomes $C\eps_Q^2\|\lambda\|_B^2$ and is absorbed.
\item Theorem~\ref{thm:reg-scalar}(b): hypothesis (ii) of Lemma~\ref{lem:CT-form} uses that $Q'^TQ'$ is block-diagonal, by (iv), and bounded by $(1+\eps_Q)^2$.
\item Theorem~\ref{thm:reg-scalar}(c): the trial field $\varphi$ built with the model transports in the local gauge satisfies $\|D\varphi\|\le C\|g\|$, as before. For the general pair, $Q'\varphi=g+r$ with $\|r\|\le C(\eps+\eps_Q)\|g\|$, so $\ip{Q'^Tg}\varphi\ge(1-C(\eps+\eps_Q))\|g\|^2$. The bound $M'_X\ge c^2$ uses $\|Q'^Tg\|=\|g\|$, which is (i).
\item Theorem~\ref{thm:reg-scalar}(d),(e): algebraic. They use (i) for the surjectivity of $Q'$, and (iv) for the facts that $Q'$ maps $\ell^2_0(X)$ to $\ell^2(X_1)$ and that $Q'^TQ'$ does not couple blocks.
\item Theorem~\ref{thm:reg-scalar}(g) and Theorem~\ref{thm:cov-scalar}(b),(d): $a(Q'^TQ'-Q_0^*Q_0)$ is pointwise admissible with constants $(C\eps_Q,0,1)$ (Lemma~\ref{lem:blocks}(c) with $C\eps$ replaced by $\eps_Q$). The pointwise bound on $H'$ uses $|Q'^Tg|_{B(y)}\le(1+\eps_Q)|g(y)|$, from the transposed bound in Definition~\ref{def:avg-parameter}(a).
\item Lemmas~\ref{lem:slice}, \ref{lem:Api-identities} and~\ref{lem:domain-identities}: (iii), with the bounds on $E$ from Definition~\ref{def:avg-parameter}(b). The real configuration $V$ enters only through $D_V$, whose transport factors are unitary.
\item Lemma~\ref{lem:local-coercive}, Step~2: $a(\bQ^T\bQ-\bQ^*\bQ)$ and the terms of $\mathcal K_2$ are admissible perturbations of size $C(\eps_Q+\eps_Q^2)$. The comparison of the scalar ingredients uses (A1)--(A5).
\item Theorem~\ref{thm:reg-coercive} and Theorem~\ref{thm:reg-gauge}(a): $\bQ^T\bQ$ is local and bounded by $(1+\eps_Q)^2$, by (ii) and (iv). This is all that the double commutators (Lemma~\ref{lem:double-comm}) and hypothesis (ii) of Lemma~\ref{lem:CT-form} use.
\item Theorem~\ref{thm:reg-gauge}(b): in the local gauge, $a(\bQ^T\bQ-\bQ^*\bQ)$ is pointwise admissible of type~$1$ with constants $(C\eps_Q,0,1)$, and $a\bQ^*\bQ$ with $(0,C,1)$. The terms $\bQ^TESu$, $DRG'(E^T\bQ u)$ and $DRG'(E^TESu)$ of $\mathcal K_2$ are of the form $\bQ^T(\text{bounded})$ or $DRG'(\text{bounded})$ with a factor $\eps_Q$; this uses the transposed bound on $E$ of Definition~\ref{def:avg-parameter}(b). Only sup bounds of the averages enter. No H\"older bound on their kernels is needed, since they contribute to $F_u$ only.
\item Lemma~\ref{lem:trial-gauge}: the model trial fields in the local gauge, built from products of $\cos^2$, with the Neumann series around the diagonally dominant matrix $T_U$. The margin is $m_\perp(\sigma_0-\sigma_+-\sigma_-)>\frac3{10}m_\perp$, with $m_\perp=2^{1-d}$, for every odd $N\ge3$. The general pair adds a local error of size $C(\eps+\eps_Q)$, which is absorbed for $\eps\le\eps_{\rm U}$ and $\eps_Q\le\theta_Q$. The slice correction uses $\|ESL\|\le C\eps_Q$.
\item Theorem~\ref{thm:reg-min}: $\bQ_\pi\bG\bQ_\pi^T$ has decay, with constants controlled by (ii) and (iv). The pointwise bound on $\bH B=\bG J$, $J=\bQ_\pi^T\boldsymbol\Gamma_\pi B$, uses three facts:
  \begin{itemize}
  \item $|\bQ^T\omega|_{B(y)}\le(1+C\eps_Q)\max|\omega|$, the maximum over the bonds $c$ with $S_c$ meeting $B(y)$ (Lemma~\ref{lem:avg}(c) for $\bQ^*$, and the transposed bound of Definition~\ref{def:avg-parameter}(a));
  \item the identity $D^T\bQ^T=Q'^TD_V^T+E^T$, from (iii), with the transposed bound on $E$ of Definition~\ref{def:avg-parameter}(b);
  \item $|Q'^Tg|\le(1+\eps_Q)|g|$.
  \end{itemize}
  The transposed sup bounds of the pair and of $E$ enter here, and also in (A5) and (A9). The pointwise bound on $\bC$ uses in addition the forward bounds from sup norms to values of $\bQ$ and of $E$ (Definition~\ref{def:avg-parameter}(a),(b)), through $\bQ_\pi\bG J=\bQ\bG J-ES\bG J$. In each case it enters only as a factor $1+C\eps_Q$, so it affects the smallness condition, not the constants.
\end{enumerate}
The explicit contour transports $\Gamma_{y,x}$ of the model pair enter only in (A3) and (A10), as stated in Remark~\ref{rem:general-averaging}. The reality in (i) is used in the G\aa rding inequality (Lemma~\ref{lem:garding}) and in the variational characterizations of $H'$ and $\bH$. No other property of the pair is used.

\emph{Step 3: the constants.} Choose $\theta_Q$, and $\eps_{\rm U}$ not larger than the threshold $\eps_0$ of Parts~II--III, below the thresholds of (A1)--(A11), of Lemma~\ref{lem:CT-stable}, and of Lemma~\ref{lem:bootstrap} (that is, $\eps_*$, depending on $d,\alpha,\gamma$). Then:
\begin{itemize}
\item Theorem~\ref{thm:reg-scalar} gives $c,C,\gamma$ depending on $d,a,\alpha,\dim\mathfrak g$, under $(1,\eps)$-regularity.
\item Lemma~\ref{lem:local-coercive} and Theorem~\ref{thm:reg-coercive} give $\frac{c_g}2$ and $c_*$ for $M\ge M_0$. In their proofs $M_0$ is fixed by $Ce^{-\gamma M_0}\le c_g/8$ and by the absorption of the terms $CM_0^{-2}$, with $C,\gamma$ from the previous item; so $M_0$ is the constant of Part~III and depends only on $d,a,\alpha,\dim\mathfrak g$. In Step~2 of the proof of Lemma~\ref{lem:local-coercive} the pair at $U^u$ is compared directly with the model pair $(Q_0,\bQ)$ at $U=1$, by Definition~\ref{def:avg-parameter}(a) and (A7); so Step~3 concerns the model pair only. For Ba\l aban's pair the pair at $U=1$ is the model pair anyway (Proposition~\ref{prop:linearized-averages}(c)).
\item Theorem~\ref{thm:reg-gauge}(a) with $X=\Tf$ follows from Lemma~\ref{lem:CT-form}, with $c_*$ and a conjugation constant $C_1$ built from the scalar decay constants. The rate is $\gamma_0=\min(\alpha_1,\frac13,c_*/(2C_1))$. It does not involve $M$.
\item Theorem~\ref{thm:reg-gauge}(b) follows from Lemma~\ref{lem:bootstrap-local}, with $\gamma'=\frac12\min(\gamma,\gamma_0)$ and a constant depending on $d,\alpha,\gamma,\gamma_0,C_\delta$. Here $C_\delta$ comes from the pointwise bounds of Theorem~\ref{thm:reg-scalar}(g). The letter $M$ in Lemma~\ref{lem:bootstrap-local} is the $L^2$ amplitude of Lemma~\ref{lem:bootstrap}, not the regularity scale. That lemma is used with its invariant inputs: the small inputs are $|u|+|\nabla u|$ on the cubes $\Lam_1(c_z)$, with weights $e^{-\gamma|y-z|}$ and the full covariant gradient ($|u|+|Du|$ for scalar fields). Only the small differential terms multiply them; the large smoothing terms use the local $L^2$ bounds from Combes--Thomas. The terms containing the pair or $E$ in (A5), (A9) and (A11) use their forward and transposed value bounds, not lattice derivatives or H\"older quotients of the averaging kernels. Every composition uses a positive rate no larger than the inverse and perturbation rates (Lemma~\ref{lem:perturbed-inverse}).
\item Theorem~\ref{thm:reg-min}: $\boldsymbol\Gamma_\pi$ has decay by Lemma~\ref{lem:CT-bounded}, with the lower bound from (A10). The sum form of the pointwise bound of $\bG$ (Theorem~\ref{thm:cov-gauge}(c), at regular backgrounds) and (A11) give (b). Here $\gamma_H$ is half the smallest of the rates above, and $C_H$ is a product of the constants above and $S_{\gamma_H}$ (Lemma~\ref{lem:decay-calculus}).
\end{itemize}
Every constant in this chain depends only on $a,\alpha,\dim\mathfrak g$. This proves (a) and (b).

\emph{Step 4: (c).} By Proposition~\ref{prop:linearized-averages}(a), (b), (d), (f), Ba\l aban's pair has properties (i), (iii) (with $E=0$ and $V=W_k$ real) and (iv). An $(M,\eps)$-regular $U_0$ has curvature at most $C_{\rm curv}\eps$ (proof of Corollary~\ref{cor:parts-apply}), so $U_0\in\mathfrak U_0(C_{\rm curv}\eps)$. Choose $c_Q$ so that $\eps\le c_Q$ implies $\max(1,C_{\rm curv})\eps\le c_3$ and the hypotheses of Lemma~\ref{lem:avg-local-gauge}. Then, in each gauge of Definition~\ref{def:avg-parameter}, Proposition~\ref{prop:linearized-averages}(e) applies on $\Lam_4(c_y)$, with $\eps$ replaced by $\max(1,C_{\rm curv})\eps$. It gives the kernel bounds with a constant that depends on $L$, by the conventions of Part~IV. Call the resulting constant $C_Q$. The kernel bounds give the block norms by the Schur test, and the sup bounds of the transposes by column sums. Hence $\eps_Q\le C_Q\eps$, and $\eps_Q\le\theta_Q$ for $\eps\le\theta_Q/C_Q$.
\end{proof}

\begin{remark}[Where $L$ and $M$ do enter]\label{rem:uniformity-limits}
\begin{enumerate}[label=(\alph*)]
\item Part~IV provides only an $L$-dependent constant $C_Q$ in Lemma~\ref{lem:uniformity}(c). The corrections $r_j$ of Lemma~\ref{lem:one-step-expansion} are carried by the averaged relative transports $\lambda_W(c_-,x)$ and by the transports $R(W(c_-,x))$, for all $x$ in the block. All contours of the families $G(c_-,x)$ start at the block centre $c_-$, so a bond near $c_-$ lies on contours to a large fraction of the $L^d$ points of the block. So the kernel of $\bQ_k(U_0)-\bQ$ may concentrate near the centres of the scale-$j$ blocks, and its sup norm, which controls the transposes, is bounded only by $O(C_L\eps)$, by the bookkeeping in the proof of Proposition~\ref{prop:Q-expansion}(c). Whether this is sharp is not settled here. Lemma~\ref{lem:uniformity} uses the transposes in (A5), (A9) and (A11). This is why \eqref{eq:H-pointwise} is stated only for $\eps$ small depending on $L$. This only affects how small the regularity parameter must be, not $C_H$ or $\gamma_H$.
\item Two results of Part~III have constants that depend on $M$. The random-walk expansion (Theorem~\ref{thm:random-walk}(a)) gives decay at a rate of order $M^{-1}\log M$, and Theorem~\ref{thm:reg-locality}(b) has the factor $e^{CM}$. Neither is used in Lemma~\ref{lem:uniformity}: there, decay comes from Combes--Thomas. Since $M\ge5L$ in Part~VI, no bound that must be independent of $L$ may use them.
\item The Dirichlet operators $\fA^{X,\kappa}$ for $X\ne\Tf$ depend on $M$ through the mass profile $m_X$. Their constants are still independent of $M\ge M_0$ (Theorem~\ref{thm:reg-coercive}), but Step~0 does not apply to them, because the operators themselves change with $M$.
\end{enumerate}
\end{remark}

Let $C_H$, $\gamma_H$ be the constants of Lemma~\ref{lem:uniformity}(b). By Lemma~\ref{lem:uniformity}(c), \eqref{eq:H-pointwise} holds at every real $(M,\eps)$-regular background $U_0$ with $M\ge M_0$ and $\eps\le\eps_{\rm B}$, with $\bH(U_0)$ the minimizer map for Ba\l aban's pair $(Q'_k(U_0),\bQ_k(U_0))$. Put
\[
C_*=\max\Big(1,\ 2dC_H\sum_{w\in\Z^d}e^{-\gamma_H|w|}\big(|w|_\infty+2\big)\Big).
\]
Then $C_*$ depends only on $a$, $\alpha$ and $G$; in particular it does not depend on $L$.


\begin{proposition}[Curvature of background fields]\label{prop:sharp-curvature}
For every $\delta>0$ there are $a_\delta\in(0,a_1]$ and an integer $n_\delta\ge n_0$, both depending on $d$, $L$, $G$, $a$, $\alpha$, $M$ and $\delta$, such that for all $k\ge0$, $n\ge n_\delta$ and $V\in\mathfrak U_k(\eps_1)$ with $\eps_1\le a_\delta$, every element of the minimal orbit satisfies $U_k(V)\in\mathfrak U_0\big((C_*+\delta)\eps_1\big)$.
\end{proposition}

\begin{proof}\hasdetail{prop:sharp-curvature}
For $k=0$, $U_0(V)=V$ and $C_*\ge1$. Let $k\ge1$ and $U=U_k(V)$. By Theorem~\ref{thm:balaban}, $U$ is $(M,\eps_0)$-regular and lies in $\mathfrak U(\eps_0)$ with $\eps_0=B_3\eps_1$, and $U$ is critical. Fix $R\ge64(M+1)$ as below and a point $y\in\T$. For $\eps_0$ small enough, depending on $R$, Steps~1--3 of the proof of Lemma~\ref{lem:improvement} apply with this $R$. This gives a real gauge transformation $w$ and, on the bonds of $\Lam_{R/4}(y)$, $U^w=e^{\ii\eta\tilde A}$ with $\tilde A=R(u)A^\sharp_*$. Here $A^\sharp_*$ is the Landau representative relative to the spliced reference $U^\sharp$, for which $\eps^\sharp+\beta^\sharp\le K_R\eps_0$.

Let $p$ be a plaquette of $\Tf$ with base point in $B(y)$. Its curvature satisfies $|F_U(p)|=|F_{e^{\ii\eta\tilde A}}(p)|\le|(\partial\tilde A)(p)|+C|\tilde A|^2$, where $\partial\tilde A(p)=\partial_\mu\tilde A_\nu-\partial_\nu\tilde A_\mu$; this follows from Lemma~\ref{lem:log-estimates}(c). By Theorem~\ref{thm:variational}(c), $A^\sharp_*=\bH^\sharp B^\sharp-\bC^\sharp J_{U^\sharp}+\mathcal R$ with $|\mathcal R|_{C^1}\le C(K_R\eps_0)^2$, where $\bH^\sharp=\bH(U^\sharp)$ and $\bC^\sharp=\bC(U^\sharp)$.
\begin{itemize}
\item In the gauge $u$, $\nabla_{U^\sharp}$ is the ordinary difference near $y$. The background $U^\sharp$ is real and $(M,\eps^\sharp)$-regular with $\eps^\sharp\le K_R\eps_0$ (Lemma~\ref{lem:splice}). We assume $K_RB_3a_\delta\le\eps_{\rm B}$ of Lemma~\ref{lem:uniformity}(c); this is one more condition on $a_\delta$, which already depends on $L$ and $M$. By \eqref{eq:H-pointwise} applied at $U^\sharp$ (Lemma~\ref{lem:uniformity}(c)),
\[
|\partial(R(u)\bH^\sharp B^\sharp)|_{B(y)}\le C_H\sum_{y'}e^{-\gamma_H|y-y'|}|B^\sharp(y')| .
\]
\item For $|y'-y|_\infty\le R/4-3$, Step~4 of that proof gives $|B^\sharp(y')|\le d(|y'-y|_\infty+2)\eps_1$. These terms contribute at most $\frac12C_*\eps_1$ to $|\partial\tilde A_\mu|$.
\item The other terms, the term $\bC^\sharp J_{U^\sharp}$ and the remainder contribute at most $CK_R\,T(R/8)\eps_0+C(K_R\eps_0)^2$, as in Step~4. Here $T$ is defined as in that proof, with $\gamma_7$ replaced by the smaller of $\gamma_H$ and the decay rate of $\bC^\sharp$.
\end{itemize}
Altogether
\[
|F_U(p)|\le C_*\eps_1+C\big(K_RT(R/8)+K_R^2\eps_0\big)B_3\eps_1+C(K_R\eps_0)^2 .
\]
Choose $R$ with $CK_RT(R/8)B_3\le\delta/2$, then $a_\delta\le a_1$ so that the remaining terms are at most $\frac\delta2\eps_1$ for $\eps_0=B_3\eps_1\le B_3a_\delta$; the bound $a_\delta\le a_1$ is needed for Theorem~\ref{thm:balaban}. Every plaquette has its base point in some block. The torus must contain the cubes $\Lam_{R+3}(y)$ without wrapping, and $n\ge n_0$ is needed for Theorem~\ref{thm:balaban}; this fixes $n_\delta=\max(n_0,4R+16)$.
\end{proof}

The proposition makes the small-field sets of consecutive steps compatible. Let $h$ and $h'$ be the curvature thresholds of the densities at steps $k$ and $k+1$, each in the units of its own step. If $V'$ has curvature at most $h'$ in the units of step $k+1$, then the curvature of $U_{k+1}(V')^{(k)}$, in the units of step $k$, is at most $L^{-2}(1+Ch')(C_*+\delta)h'$, by Proposition~\ref{prop:avg-regular}. In four dimensions the change of units does not change the coupling, so the thresholds of consecutive steps are comparable. We use only $h'\le2h$. For thresholds $h=p(g_k)g_k$ with $p$ decreasing and $g\mapsto p(g)g$ increasing, as in Section~\ref{sec:flow}, this holds whenever $g_{k+1}\le2g_k$. With $h'\le2h$, $\delta\le1$ and $Ch'\le1$, the bound is at most $4(C_*+1)L^{-2}h$, and this is at most $\frac12h$ as soon as
\[
L^2\ge8(C_*+1).
\]
Since $C_*$ does not depend on $L$, this is a condition on $L$ alone. In the flow the thresholds are $h_k=p(g_k)g_k$ (Section~\ref{subsec:flow-parameters}), and $g_{k+1}\le2g_k$, hence $h_{k+1}\le2h_k$, holds at every step (Theorem~\ref{thm:gen-step}(c), Theorem~\ref{thm:sf-flow}(e)).

\subsection{The background of a step}\label{ss:31.3}

\begin{lemma}[Composition of background fields]\label{lem:composition}
Let $V'$ be a real configuration on $\T'$ whose curvature, in the units of step $k+1$, is at most $\eps'_1\le\min\big(a_{\delta=1},\,a_1/(2C_*+2)\big)$, with $a_\delta$ from Proposition~\ref{prop:sharp-curvature}, and $(C_*+1)\eps'_1\le c_1$. Let $U'=U_{k+1}(V')$ be an element of its minimal orbit, and $V^0=U'^{(k)}$. Then:
\begin{enumerate}[label=(\alph*)]
\item $V^0\in\mathfrak U_k(\eps_1)$, in the units of step $k$, with $\eps_1=2(C_*+1)L^{-2}\eps'_1\le a_1$, and $\mathfrak a'(V^0)=V'$;
\item $U'$ belongs to the minimal orbit of the $k$-step problem for $V^0$: $U_k(V^0)=U'$;
\item for every real $V$ with $\mathfrak a'(V)=V'$ and $V\in\mathfrak U_k(a'_1)$, $a'_1=\min(a_1,\theta_0L^{-3}B_3^{-1})$, we have $S(U_k(V))\ge S(U')$. Equality holds iff $V=(V^0)^u$ for a real $u$ with $u|_{\T'}=1$.
\end{enumerate}
\end{lemma}

\begin{proof}\hasdetail{lem:composition}
(a) By Proposition~\ref{prop:sharp-curvature} (with $\delta=1$), $U'$ has curvature at most $(C_*+1)\eps'_1$ in the units of step $k+1$. By Proposition~\ref{prop:avg-regular} its $k$-fold average has curvature at most twice this, in the same units, which is the claim in the units of step $k$. Moreover $\mathfrak a'(V^0)=U'^{(k+1)}=V'$.

(b) Near $U'$, the constraint manifold $\{U^{(k)}=V^0\}$ is a submanifold of $\{U^{(k+1)}=V'\}$, since $U^{(k)}=V^0$ implies $U^{(k+1)}=\mathfrak a'(V^0)=V'$. As sets, $\mathfrak B_k(V^0)\subset\mathfrak B_{k+1}(V')$ is false: the cap $U\in\mathfrak U_0(c_1)$ in the definition of $\mathfrak B$ is weaker in the units of step $k$. Only the local inclusion is used. The configuration $U'$ is a critical point of $S$ on the manifold $\{U^{(k+1)}=V'\}$ near $U'$ and lies on the submanifold $\{U^{(k)}=V^0\}$. So it is critical on $\mathfrak B_k(V^0)$. In the units of step $k$ its curvature and current are smaller than in the units of step $k+1$, so $U'\in\mathfrak U(\theta_0)$ in the units of step $k$. By Theorem~\ref{thm:balaban}(a) at step $k$, $U'$ lies in the unique critical orbit, which is the minimal orbit.

(c) By Theorem~\ref{thm:balaban}(b), $U_k(V)\in\mathfrak U(B_3a'_1)$ in the units of step $k$. In the units of step $k+1$ this gives curvature and current at most $L^3B_3a'_1\le\theta_0$. Hence $U_k(V)\in\mathfrak B_{k+1}(V')\cap\mathfrak U(\theta_0)$, and $S(U_k(V))\ge S(U')$. If equality holds, $U_k(V)$ is a minimizer of the $(k+1)$-step problem, so $U_k(V)=U'^{u}$ with $u|_{\T'}=1$. Then $V=U_k(V)^{(k)}=(V^0)^{u|_\T}$. The converse follows from gauge invariance.
\end{proof}

So $V^0$ is the minimizer of the coarse action $V\mapsto S(U_k(V))$ on the one-step constraint set, and the background of the $(k+1)$-step problem is the background of $V^0$. We call $V^0$ the \emph{base configuration} of the step. It depends on the choice of $U'$ in its orbit only through a gauge transformation on $\T$ that is trivial on $\T'$.

In the units of step $k$, $U'$ is $(M,B_3\eps'_1)$-regular and lies in $\mathfrak U(B_3\eps'_1)$. Indeed, passing from the units of step $k+1$ to those of step $k$ multiplies Lie fields by $L^{-1}$, curvatures by $L^{-2}$ and currents by $L^{-3}$, and the cubes of Definition~\ref{def:regular} become larger, which only helps (Lemma~\ref{lem:regular-monotone}). Throughout this part we write
\[
\eps:=B_3\eps'_1
\]
for the regularity parameter of $U'$ as a reference configuration for the $k$-step problem.

\subsection{Lattice symmetry}\label{ss:31.4}

Let $\mathfrak S'=\mathfrak S_{k+1}$ be the group of the maps $x\mapsto Rx+v$ of $\Tf$ with $R\in B_4$ and $v\in\T'$ (Section~\ref{sec:averaging}), in the units of step $k$. It is generated by the translations of $\T'$ and by $B_4$ about the points of $\T'$, and $\mathfrak S'\subset\mathfrak S_k$. Its elements map $\T'$, $\T$ and $\Tf$ onto themselves, and $B'(y)$ onto $B'(ry)$. They act on configurations by $(rV)(b)=V(r^{-1}b)$, on scalar fields by $(r\lambda)(x)=\lambda(r^{-1}x)$, on functions of configurations by $(r\rho)(V)=\rho(r^{-1}V)$, and on Lie fields relative to a configuration $U$ by $r_U$ (Proposition~\ref{prop:linearized-averages}(g)). Let $\mathcal G'$ be the group of real gauge transformations on $\T$ that are $1$ on $\T'$. The following proposition collects the lattice symmetries of the objects constructed so far.

\begin{proposition}[Lattice symmetry]\label{prop:symmetry}
\begin{enumerate}[label=(\alph*)]
\item \emph{(The map $T$.)} For $r\in\mathfrak S'$ and every density $\rho$ on configurations on $\T$ that vanishes outside $\operatorname{dom}\mathfrak a'$, $T(r\rho)=r\,T(\rho)$. If $\rho$ is gauge invariant, so is $r\rho$.
\item \emph{(Background fields.)} Let $r\in\mathfrak S_k$ and $V\in\mathfrak U_k(\eps_1)$ with $\eps_1\le a_1$. Then $U_k(rV)$ and $rU_k(V)$ lie in the same $\mathfrak G_0$-orbit, and $S(U_k(rV))=S(U_k(V))$. Let $r\in\mathfrak S'$ and let $V'$ be as in Lemma~\ref{lem:composition}. Then $rV'$ satisfies the same hypotheses, and the base configurations satisfy $V^0(rV')=(rV^0(V'))^u$ for some $u\in\mathcal G'$.
\item \emph{(Averages and operators.)} Let $U_0$ be real and $(M,\eps)$-regular with $\eps\le\eps_{\rm B}$ (Lemma~\ref{lem:uniformity}(c)), and let $r\in\mathfrak S_k$. Then $rU_0$ is $(M,\eps)$-regular, $J_{rU_0}=r_{U_0}J_{U_0}$ and $W_k(rU_0)=rW$, where $W=W_k(U_0)$. Moreover:
  \begin{enumerate}[label=(\roman*)]
  \item $Q'_k(rU_0)\,r=r\,Q'_k(U_0)$, $\bQ_k(rU_0)\,r_{U_0}=r_W\,\bQ_k(U_0)$, and $D_{rU_0}r=r_{U_0}D_{U_0}$ on scalar fields;
  \item $K(rU_0)\,r=r\,K(U_0)$ for the scalar operators $K=\fA',G',\Gamma',M',H',C',P_\cV,R$ of Parts~II and~III on the torus; and, for the operators $S(U_0)=G'RD^T$ and $\Pi(U_0)=I-DS(U_0)$ of Section~\ref{sec:cov-gauge}, $S(rU_0)\,r_{U_0}=r\,S(U_0)$ and $\Pi(rU_0)\,r_{U_0}=r_{U_0}\,\Pi(U_0)$;
  \item $\Phi_{rU_0}(r_{U_0}A)=\Phi_{U_0}(A)$ and $\ell\,\mathcal Q(e^{\ii\eta\,r_{U_0}A};rU_0)=r_W\,\ell\,\mathcal Q(e^{\ii\eta A};U_0)$, wherever defined; so the derivatives of all orders of these maps at $r_{U_0}A$ and at $A$ are conjugate by $r_{U_0}$ and $r_W$;
  \item let $\mathfrak B=D^TD+W'$ be the operator of the form $q_U$ in Section~\ref{sec:cov-gauge}. If $\mathfrak B(rU_0)\,r_{U_0}=r_{U_0}\,\mathfrak B(U_0)$, then the operators $\fA_\pi$, $\bG$, $\boldsymbol\Gamma_\pi$, $\bH$ and $\bC$ built with $\mathfrak B$ in Sections~\ref{sec:cov-gauge} and~\ref{sec:cov-min} are covariant in the same way; for instance $\bH(rU_0)\,r_W=r_{U_0}\bH(U_0)$ and $\bC(rU_0)\,r_{U_0}=r_{U_0}\bC(U_0)$. For $W'=0$ this holds if $r(x)=Rx+v$ with a permutation matrix $R$, and for all $r\in\mathfrak S_k$ if $U_0=1$.
  \end{enumerate}
  Finally, for $r\in\mathfrak S'$ and a configuration $Y$ on $\T$ as in Definition~\ref{def:avg-transport}, for the lattices $\T$ and $\T'$: $(\mathfrak G_{rY}r_YB)(y,x)=(\mathfrak G_YB)(r^{-1}y,r^{-1}x)$.
\item \emph{(Landau representatives.)} Let $V$ be as in Theorem~\ref{thm:balaban}, $U_0$ an element of its minimal orbit, and $r\in\mathfrak S_k$. Then $rU_0$ is an element of the minimal orbit of $rV$. In the notation of Corollary~\ref{cor:analytic}, $A_*(r_VB;rU_0)=r_{U_0}A_*(B;U_0)$, and, for real $B$, $\omega(r_VB;rU_0)=r_V\,\omega(B;U_0)$.
\end{enumerate}
\end{proposition}

\begin{proof}\hasdetail{prop:symmetry}
(a) Substitute $V=rY$ in \eqref{eq:rg-map}. The map $Y\mapsto rY$ permutes the bond variables and inverts those on the bonds $b$ for which $r^{-1}b$ is negatively oriented. The Haar measure is invariant under inversion, so $\prod_bdV(b)=\prod_bdY(b)$, and the density becomes $(r\rho)(rY)=\rho(Y)$. By Lemma~\ref{lem:one-step-basic}(d), $\mathfrak a'(rY)$ is defined if and only if $\mathfrak a'(Y)$ is, so $r\rho$ vanishes outside $\operatorname{dom}\mathfrak a'$, and $\mathfrak a'(rY)=r\,\mathfrak a'(Y)$. Put $X=\mathfrak a'(Y)$ and $\tilde V'=r^{-1}V'$. Then $\mathfrak a'(rY)(c)\,V'(c)^{-1}=X(r^{-1}c)\,\tilde V'(r^{-1}c)^{-1}$. As $c$ runs over the positively oriented bonds of $\T'$, $r^{-1}c$ runs over all bonds of $\T'$ once, up to orientation. For a reversed bond, $X(-c)\tilde V'(-c)^{-1}=X(c)^{-1}\tilde V'(c)$ is conjugate to $(X(c)\tilde V'(c)^{-1})^{-1}$, and $\delta$ is invariant under conjugation and inversion. Hence the product of $\delta$-functions is $\prod_c\delta(X(c)\tilde V'(c)^{-1})$, and $(T(r\rho))(V')=(T\rho)(r^{-1}V')$. Finally $r(V^u)=(rV)^{ru}$, so $r\rho$ is gauge invariant if $\rho$ is.

(b) The first two claims are Theorem~\ref{thm:balaban}(c). The curvature of $rV'$ equals that of $V'$, since $\mathfrak U_{k+1}(\eps'_1)$ is invariant under $\mathfrak S'$. In the units of step $k+1$, $\mathfrak S'$ is the group of the lattice symmetries of the $(k+1)$-step structure. By Theorem~\ref{thm:balaban}(c) for the $(k+1)$-step problem, $U_{k+1}(rV')=(rU')^u$ with $U'=U_{k+1}(V')$ and a real $u$ on $\Tf$ with $u|_{\T'}=1$. By the covariance of $\mathfrak a^{(k)}$ under real gauge transformations and under $\mathfrak S_k\supset\mathfrak S'$ (Lemma~\ref{lem:one-step-basic}),
\[
V^0(rV')=\big((rU')^u\big)^{(k)}=\big(r\,U'^{(k)}\big)^{u|_\T}=\big(rV^0(V')\big)^{u|_\T},
\]
and $u|_\T\in\mathcal G'$.

(c) The regularity of $rU_0$ and the identity for $W_k$ are proved in the proof of Theorem~\ref{thm:variational}(h) (Section~\ref{subsec:var-symmetry}), and the identity for $J$ is Lemma~\ref{lem:current}(b). An $(M,\eps)$-regular $U_0$ with $\eps\le\eps_{\rm B}$ lies in $\mathfrak U_0(c_3)$ (Step~4 of the proof of Lemma~\ref{lem:uniformity}), so Proposition~\ref{prop:linearized-averages} applies. (i) The first two identities are Proposition~\ref{prop:linearized-averages}(g), and the third is shown in Section~\ref{subsec:var-symmetry}. (ii) The maps $r$ and $r_{U_0}$ preserve the bilinear pairings, so transposes are covariant as well. The operators in (ii) are built from $D$ on scalar fields, $D^T$, $Q'_k$ and $Q'^T_k$ by sums, products, transposes and inverses (Sections~\ref{sec:cov-scalar} and~\ref{sec:cov-gauge}), so (i) gives (ii). (iii) The first identity follows from $e^{\ii\eta\,r_{U_0}A}\,rU_0=r(e^{\ii\eta A}U_0)$ and $S\circ r=S$; the second is Proposition~\ref{prop:linearized-averages}(g), together with $\ell((rX)(c))=(r_W\ell X)(c)$ for relative configurations at scale $k$. Differentiating, and using that $r_{U_0}$ and $r_W$ preserve the pairings, gives the claim about derivatives. (iv) By definition, $\fA_\pi=\Pi^T\mathfrak B\Pi+DRD^T+a\Pi^T\bQ^T\bQ\Pi$, $\bG=\fA_\pi^{-1}$, $\bQ_\pi=\bQ\Pi$, $\boldsymbol\Gamma_\pi=(\bQ_\pi\bG\bQ_\pi^T)^{-1}$, $\bH=\bG\bQ_\pi^T\boldsymbol\Gamma_\pi$ and $\bC=\bG-DG'RG'D^T-\bH\bQ_\pi\bG$. With (i) and (ii), each of them is covariant if $\mathfrak B$ is. Let $W'=0$. If $R$ is a permutation matrix, $e_\mu\mapsto e_{\sigma(\mu)}$, then $r$ maps positively oriented bonds onto positively oriented bonds, $r_{U_0}$ only relabels the bonds, and directly from the definition of $D$ on gauge fields,
\[
\big(D_{rU_0}\,r_{U_0}A\big)_{\sigma(\mu)\sigma(\nu)}(rx)=\big(D_{U_0}A\big)_{\mu\nu}(x).
\]
Hence $\|D_{rU_0}r_{U_0}A\|=\|D_{U_0}A\|$, and polarization gives $\mathfrak B(rU_0)\,r_{U_0}=r_{U_0}\,\mathfrak B(U_0)$. At $U_0=1$, $D=d$ and $(r_1A)(b)=A(r^{-1}b)$ with $A(-b)=-A(b)$. Then $(dA)(p)$ is the sum of $A$ around $\partial p$, which does not depend on the corner at which $p$ is based, and $\|d\,r_1A\|=\|dA\|$ for all $r\in\mathfrak S_k$. The last claim is Lemma~\ref{lem:avg-transport}(d), applied to the lattices $\T$ and $\T'$.

(d) By Theorem~\ref{thm:balaban}(c), $rU_0$ is an element of the minimal orbit of $rV$. Since $W_k(U_0)=V$, Theorem~\ref{thm:variational}(h) gives the rest; see also Corollary~\ref{cor:analytic}.
\end{proof}

\begin{remark}[Reflections, and what the step uses]\label{rem:symmetry-reflections}
At curved backgrounds and for $W'=0$, the hypothesis of part~(c)(iv) is verified for the translations and the permutations of the axes, but not for the reflections. The covariant derivative of a Lie field on a plaquette is linearized at its base corner: $(D_UA)_{\mu\nu}(x)=D_\mu A_\nu(x)-D_\nu A_\mu(x)$, with transports from $x$. A reflection moves the base corner of a plaquette to another corner of its image. Let $r$ be the reflection in the direction $\mu$ about a point of $\T$, let $\nu\ne\mu$, $z=r^{-1}x$, $z'=z-\eta e_\mu$, and let $p=p_{\mu\nu}(z')$. A direct computation from the definitions gives
\[
\big(D_{rU}\,r_UA\big)_{\mu\nu}(x)=-R(U_\mu(z'))^{-1}\Big[(D_UA)_{\mu\nu}(z')-\eta^{-1}\big(R(U(\partial p))-1\big)R(U_\nu(z'))A_\mu(z'+\eta e_\nu)\Big].
\]
So $D_{rU}r_U$ and the transported $D_U$ differ by a local operator of norm at most $C\eta\max_p|F_U(p)|$, which vanishes for flat $U$. Since $\|D\|\le C\eta^{-1}$, the operator $D^TD$ at a background of curvature $\eps$ is covariant under reflections up to a perturbation of norm $C\eps$. For $\bH$, $\bC$ and the other operators of~(c)(iv) at curved backgrounds we therefore claim covariance only under the translations and the permutations of the axes. They are tools in the analysis. The quantities that are defined without a choice of gauge are covariant under all of $\mathfrak S'$: the map $T$, by~(a); the background fields and the coarse action $V\mapsto S(U_k(V))$, by~(b); and the derivatives of the action and of $\mathcal Q$, the averaging pair, the scalar operators and the averaged relative transports, by~(c). The covariance of the small-field step, $T_{\rm sf}(r\rho)=r\,T_{\rm sf}(\rho)$ for $r\in\mathfrak S'$, which also involves the gauge-fixing function and the cutoffs, is part~(iii) of Lemma~\ref{lem:gf-step-invariance}. The proof devices of Parts~II--VI, namely the model pair of Section~\ref{sec:background}, the contours $\Gamma^{\rm co}$ of Section~\ref{sec:local-gauges}, the spliced configurations and the trivial extension, enter only proofs. They are not covariant, and no statement above depends on them.
\end{remark}
